\chapter{The Fundamental Estimate on Curves}
\label{chap:curve-estimate}

We now prove the one-dimensional estimate from which the global
direct-image theorem will follow.  Cohomology in degree \(i\) raises an
upper weight bound by at most \(i\).  The proof uses the real-sheaf
theorem, Fourier inversion, Plancherel, the determinant-weight
\(L^2\)-criterion, local duality, and elementary geometric reductions.

\section{Statement and the intermediate-extension form}

\begin{theorem}[Curve estimate]
\label{thm:curve-estimate}
Let \(C_0\) be a smooth geometrically irreducible curve over \(k_0\),
not necessarily proper, and let \(\cF_0\) be a lisse Weil sheaf
punctually \(\iota\)-pure of weight \(w\).  Then
\[
  H_c^i(C,\cF)
\]
is \(\iota\)-mixed of weights at most \(w+i\) for every \(i\).
\end{theorem}

Let \(\bar C_0\) be the smooth projective completion and
\(j_0:C_0\hookrightarrow\bar C_0\).

\begin{theorem}[Intersection-cohomology form]
\label{thm:curve-intersection-purity}
Under the hypotheses of \cref{thm:curve-estimate},
\[
 H^i(\bar C,j_*\cF)
\]
is punctually \(\iota\)-pure of weight \(w+i\) whenever it is nonzero.
\end{theorem}

\section{Boundary weights and compact support}

The needed boundary assertion is Deligne's semicontinuity theorem.  A
local-monodromy filtration alone does not prove it: that would silently
import a weight--monodromy statement.  The following proof instead uses
an \(L\)-function and a tensor-power limit.

\begin{foundation}
\label{found:boundary-semicontinuity}
\textbf{Semicontinuity at the boundary.}
Let \(j_0:U_0\hookrightarrow C_0\) be a dense open immersion of smooth
geometrically irreducible curves, and let \(\cF_0\) be lisse and pure of
weight \(w\) on \(U_0\).  Then every Frobenius eigenvalue on a boundary
stalk of \(j_{0*}\cF_0\) has weight at most \(w\).

\smallskip
\emph{Proof.}
Fix \(s_0\in C_0-U_0\), of degree \(e\), and an eigenvalue
\(\alpha\) on
\[
  (j_*\cF)_{\bar s}=V^{I_s}.
\]
Delete a closed point of \(U_0\), if necessary, so that \(C_0\) is
affine; this changes neither the boundary stalk nor the assertion.  For
\(m\ge1\), put
\[
  \cG_{m,0}=j_{0*}(\cF_0^{\otimes m}).
\]
The Frobenius-equivariant injection
\[
  (V^{I_s})^{\otimes m}\hookrightarrow(V^{\otimes m})^{I_s}
\]
shows that \(\alpha^m\) occurs on
\((\cG_m)_{\bar s}\).

The adjunction
\(\cG_{m,0}\to j_{0*}j_0^*\cG_{m,0}\) is an isomorphism, so
\(\cG_m\) has no section supported on the boundary.  On the affine
curve this gives \(H_c^0(C,\cG_m)=0\).  Removing finitely many points
does not change \(H_c^2\), and top-degree duality gives
\[
 H_c^2(C,\cG_m)
 \simeq(V^{\otimes m})_{\pi_1(U)}(-1).
\]
The coinvariants are a geometrically constant quotient of the pure
sheaf \(\cF_0^{\otimes m}\), so the eigenvalues on this group have
weights at most \(mw+2\).  Hence the trace formula
\[
 L(C_0,\cG_{m,0},T)
 =
 \frac{\det(1-TF\mid H_c^1(C,\cG_m))}
      {\det(1-TF\mid H_c^2(C,\cG_m))}
\tag{11.1}\label{eq:boundary-trace-ratio}
\]
has no pole for
\[
  |T|<q^{-mw/2-1}.
\tag{11.2}\label{eq:boundary-pole-disk}
\]

Purity on \(U_0\), the logarithmic derivative of its Euler product, and
\(\#U_0(\mathbf F_{q^n})=O(q^n)\) show that
\(L(U_0,\cF_0^{\otimes m},T)\) converges absolutely and is nonzero in
the same disk.  Euler products factor as
\[
 L(C_0,\cG_{m,0},T)
 =
 L(U_0,\cF_0^{\otimes m},T)
 \prod_{t_0\in C_0-U_0}
 \det\!\left(
  1-T^{\deg(t_0)}F_{t_0}\mid(\cG_m)_{\bar t}
 \right)^{-1}.
\tag{11.3}\label{eq:boundary-euler-factorization}
\]
After division by the nonvanishing first factor, the finite boundary
product has no pole in \cref{eq:boundary-pole-disk}.  It has a
denominator but no numerator, so its individual local factors cannot
cancel one another.  Since \(\alpha^m\) occurs in the factor at \(s_0\),
\[
 |\iota(\alpha)|^{m/e}\le q^{mw/2+1}.
\]
Thus
\[
 |\iota(\alpha)|\le q^{e(w/2+1/m)}.
\]
Letting \(m\to\infty\) gives
\(|\iota(\alpha)|\le q^{ew/2}\).

This proof uses only the trace formula, top-degree duality, elementary
Euler-product convergence, and purity on the open; it uses neither
weight--monodromy nor the global direct-image theorem.  See
\cite[Proposition~2.3.0.1]{StanfordL19} and
\cite[Chapter~I, Section~2]{KiehlWeissauer}.
\end{foundation}

Let
\(\cQ_0=j_{0*}\cF_0/j_{0!}\cF_0\).  It is punctual of weights at most
\(w\), and
\[
 0\longrightarrow j_!\cF\longrightarrow j_*\cF
 \longrightarrow\cQ\longrightarrow0.
\tag{11.4}\label{eq:jshriek-jstar}
\]
The long exact sequence contains
\[
 0\to H_c^0(C,\cF)\to H^0(\bar C,j_*\cF)
 \to H^0(\bar C,\cQ)\to H_c^1(C,\cF)
 \to H^1(\bar C,j_*\cF)\to0
\tag{11.5}\label{eq:boundary-long-exact}
\]
and
\[
 H_c^2(C,\cF)\simeq H^2(\bar C,j_*\cF).
\]
Thus \cref{thm:curve-intersection-purity} implies
\cref{thm:curve-estimate}.

\section{Reduction to an open subset of the affine line}

Finite extension of \(k_0\) does not change normalized weights.  Delete
finitely many points so that \(C_0\) is affine.  Noether normalization
gives a finite generically \'etale map to \(\A^1\); after deleting the
branch locus it is finite \'etale.  Finite pushforward preserves purity:
Frobenius permutes the stalks above a point, and each orbit block is
obtained from the Frobenius polynomial upstairs by a power substitution.
Extending the finite map to projective completions gives
\[
 R\Gamma(\bar C,j_*\cF)
 \simeq R\Gamma(\PP^1,j_*g_*\cF).
\]
We may therefore replace the pair by an open
\(U_0\subset\PP^1_{k_0}\) and a pure lisse sheaf on it.

Pass to a Jordan--Holder constituent and, after a finite base extension,
assume it geometrically irreducible.  A geometrically constant
constituent has \(H^1(\PP^1,j_*\cF)=0\), so assume geometric
nonconstancy.  Choose \(s\in U_0(k_0)\), remove it, and send it to
\(\infty\).  Because the original sheaf extended across \(s\), this does
not change \(j_*\cF\), and the resulting sheaf on
\[
 U_0\subset\A^1_{k_0}
\]
is unramified at infinity.  We have reduced to a geometrically
irreducible, geometrically nonconstant sheaf unramified at infinity.

\section{The Fourier transform is a mixed sheaf}

Let \(j:U\hookrightarrow\A^1\), put
\[
 \cK_0=j_!\cF_0,\qquad
 \widehat{\cF}_0=\FT_\psi(\cK_0).
\]

\begin{lemma}
\label{lem:twisted-cohom-concentration}
For every geometric \(a\in\A^1(k)\),
\[
 H_c^i(U,\cF\otimes\cL_{\psi_a})=0\qquad(i\ne1).
\]
Consequently \(\widehat{\cF}_0\) is an ordinary constructible Weil
sheaf in degree \(0\).
\end{lemma}

\begin{proof}
Compactly supported \(H^0\) of a lisse sheaf on the connected affine
curve vanishes.  Top-degree duality identifies \(H_c^2\) with the
geometric coinvariants, twisted by \((-1)\).  Since
\(\cF\otimes\cL_{\psi_a}\) is geometrically irreducible, nonzero
coinvariants would make it geometrically constant.  For \(a=0\) this
contradicts the nonconstancy of \(\cF\).  For \(a\ne0\), it contradicts
the fact that \(\cF\) is unramified at infinity whereas
\(\cL_{\psi_a}\) has nontrivial wild inertia there.  The stalk formula
\cref{cor:fourier-cohomology-sheaves} proves the last assertion.
\end{proof}

\begin{lemma}
\label{lem:fourier-mixed-before-estimate}
The sheaf \(\widehat{\cF}_0\) is \(\iota\)-mixed.
\end{lemma}

\begin{proof}
Let \(\cF_0^c\) be the conjugate summand supplied by
\cref{lem:real-envelope}.  Apply the preceding construction to
\((\cF_0,\psi)\) and \((\cF_0^c,\psi^{-1})\).  Their Fourier trace
functions are complex conjugates, so
\[
 \FT_\psi(j_!\cF_0)\oplus
 \FT_{\psi^{-1}}(j_!\cF_0^c)
\]
is \(\iota\)-real.  It is mixed by
\cref{thm:real-sheaf-mixed}, and so is its direct summand
\(\widehat{\cF}_0\).
\end{proof}

\begin{lemma}
\label{lem:fourier-no-hc0}
\[
 H_c^0(\A^1,\widehat{\cF})=0.
\]
In particular, the semisimplification of \(\widehat{\cF}\) has no
punctual constituent.
\end{lemma}

\begin{proof}
The inverse-transform stalk formula and Fourier inversion give
\[
 H_c^0(\A^1,\widehat{\cF})
 \simeq
 \cH^{-1}\!\left(
  \FT_{\psi^{-1}}(\widehat{\cF})
 \right)_0
 \simeq
 \cH^{-1}\bigl((j_!\cF)(-1)\bigr)_0=0.
\]
A punctual direct summand would have a nonzero compactly supported
section.
\end{proof}

\section{Plancherel gives the upper bound}

The sheaf \(\cK_0=j_!\cF_0\) is pure of weight \(w\), has no punctual
constituent, and has \(H_c^0=0\).  Thus
\cref{thm:l2-weight-criterion} gives
\[
 \limsup_n\frac1n\log_q
 \|\iota t_{\cK}^{(n)}\|_{2,n}^2=w+1.
\tag{11.6}\label{eq:K-l2-growth}
\]
Plancherel gives
\[
 \|\iota t_{\widehat{\cF}}^{(n)}\|_{2,n}^2
 =q^n\|\iota t_{\cK}^{(n)}\|_{2,n}^2.
\]
The left exponent is therefore \(w+2\).  The Fourier sheaf is mixed and
has neither \(H_c^0\) nor punctual constituents, so applying the
\(L^2\)-criterion again gives
\[
  \widehat w_{\max}=w+1.
\tag{11.7}\label{eq:fourier-weight-shift}
\]
At \(a=0\), the stalk formula identifies
\[
 \widehat{\cF}_{\bar0}\simeq H_c^1(U,\cF).
\]

\begin{proposition}[The fundamental upper estimate]
\label{prop:fundamental-curve-upper}
Every eigenvalue on \(H_c^1(U,\cF)\) has weight at most \(w+1\).  The
same holds on its quotient \(H^1(\PP^1,j_*\cF)\).
\end{proposition}

\section{Duality gives equality}

We have
\[
 H^0(\bar C,j_*\cF)=H^0(C,\cF).
\]
Its nonzero constituents are geometrically constant quotients of
\(\cF\), hence have weight \(w\).  Local duality gives
\[
 H^2(\bar C,j_*\cF)
 \simeq H^0(\bar C,j_*\cF^\vee)^\vee(-1),
\]
so its nonzero eigenvalues have weight \(w+2\).

In degree one, local duality gives a perfect pairing
\[
 H^1(\bar C,j_*\cF)\otimes
 H^1(\bar C,j_*\cF^\vee)
 \longrightarrow\Qbarl(-1).
\tag{11.8}\label{eq:curve-local-duality}
\]
If \(\alpha\) occurs in the first factor, \(q/\alpha\) occurs in the
second.  Apply \cref{prop:fundamental-curve-upper} to the dual sheaf,
which has weight \(-w\):
\[
 |\iota(q/\alpha)|\le q^{(-w+1)/2}.
\]
Thus
\[
 |\iota(\alpha)|\ge q^{(w+1)/2}.
\]
The upper estimate for \(\cF\) gives the reverse inequality.  This
proves \cref{thm:curve-intersection-purity}; the boundary exact sequence
then proves \cref{thm:curve-estimate}.

\begin{warning}
\label{warn:curve-proof-noncircular}
The boundary estimate came from an \(L\)-function tensor-power limit,
not from weight--monodromy.  The lower middle-degree estimate came from
the already proved upper estimate applied to the dual, not from
smooth-proper purity.
\end{warning}

\begin{formalization}
\label{form:curve-estimate-interface}
\textbf{Formalization decomposition.}
\begin{enumerate}
\item Boundary stalks as inertia invariants and Deligne's
\(L\)-function/tensor-power semicontinuity proof.
\item The \(j_!\)-to-\(j_*\) exact sequence.
\item Finite-map reduction to \(\PP^1\).
\item Vanishing for every Artin--Schreier twist.
\item Fourier transform as an ordinary mixed sheaf with \(H_c^0=0\).
\item Plancherel plus the determinant-weight \(L^2\)-criterion.
\item Local duality for the reverse bound.
\item The terminal implication
\[
 \operatorname{Pure}(\cF,w)
 \Longrightarrow
 \operatorname{WeightsLE}(H_c^i(C,\cF),w+i).
\]
\end{enumerate}
No step imports the global \(Rf_!\) theorem.
\end{formalization}
